\section{An incidence piece and its probability laws}\label{sec:incidence}

The input to this section is a sequence of planar configurations with
\(t=n^{1/3}\to\infty\) and at least \(t^{4-o(1)}\) ordered unit pairs.
We construct an unbalanced incidence graph whose degrees are controlled
on both sides. Its edge law has a uniform expansion estimate, including
for arbitrarily small sets. We then prove two consequences of that
estimate and show that the endpoints of a random two-edge path have
mutual information \(o(\log t)\).

Such an input sequence exists if \Cref{thm:main} fails. Indeed, for each
\(j\ge2\), failure of a bound with exponent \(4/3-1/j\) allows the
constant in that bound to be chosen large enough to force a counterexample
with \(n\ge j\). Its ordered unit-pair count is at least
\(n^{4/3-1/j}=t^{4-3/j}\). Passing to an increasing sequence of sizes
gives the claimed input. We use separate copies of the configuration as
points and as centers of unit circles, so its ordered unit pairs are
exactly the incidences between these two copies.

\subsection{Preserving incidences under algebraic specialization}

Later arguments use number-field absolute values. The following
specialization permits this without losing any of the unit pairs.

\begin{lemma}\label{inc:specialization}
Let \(x_1,\ldots,x_n\) be distinct points of \(\R^2\), and let
\(\mathcal E\) be any collection of pairs \((i,j)\) for which
\(\norm{x_i-x_j}_2=1\). There are distinct points
\(x'_1,\ldots,x'_n\) with real algebraic coordinates, arbitrarily close
to the respective original points, such that
\(\norm{x'_i-x'_j}_2=1\) for every \((i,j)\in\mathcal E\).
\end{lemma}

\begin{proof}
Fix a neighborhood of each original point and choose a rational open
rectangle containing that point inside its neighborhood. For variables
\(X_i=(X_{i,1},X_{i,2})\), impose membership in these rectangles,
the equations
\[
 (X_{i,1}-X_{j,1})^2+(X_{i,2}-X_{j,2})^2=1
       \qquad ((i,j)\in\mathcal E),
\]
and the strict inequalities
\[
 (X_{i,1}-X_{j,1})^2+(X_{i,2}-X_{j,2})^2>0
       \qquad (i\ne j).
\]
This finite system has rational coefficients and a real solution, namely
the original configuration. The transfer principle for real closed fields
\cite[Theorems~3.1 and~4.1]{RCFTransfer} gives a solution in the field
of real algebraic numbers. The equations preserve the required unit pairs,
the inequalities preserve distinctness, and the rectangles give the
prescribed closeness.
\end{proof}

We henceforth apply this lemma separately to each configuration. Extra
unit pairs created by the specialization are harmless; all subsequent
graphs may retain only the incidences under consideration.

\subsection{Cuttings and the incidence bound}

For finite sets \(\mathcal P\) of points and \(\mathcal Q\) of distinct
unit-circle centers, let \(I(\mathcal P,\mathcal Q)\) denote the number
of incidences. Two distinct unit circles meet in at most two points.
Also, two distinct points lie on at most two unit circles: their possible
centers are intersections of the two unit circles centered at the points.
These two facts provide the pair-counting bounds used below.
The cutting proof uses the random-sampling principle for regions with
bounded defining data developed by Clarkson \cite[Section~4]{Clarkson1987}.
We give the circle boundary conventions explicitly so that every retained
point belongs to its assigned open region.

\begin{lemma}\label{inc:cutting}
Let \(\mathcal P\) be a finite planar point set and let \(\mathcal Q\)
be a nonempty set of \(l\) distinct unit circles. For every \(R\ge1\),
sample
\(s=\lceil C R\log(2+l)\rceil\) circles independently and uniformly
with replacement, where \(C\) is a sufficiently large absolute constant.
With probability at least \(9/10\), after deleting the points on sampled
circles, all remaining points can be assigned to at most \(C'(1+s^2)\)
open regions, each of which intersects at most \(l/R\) of the original
circles. Each assigned point belongs to its assigned region.
\end{lemma}

\begin{proof}
Choose an auxiliary horizontal direction before taking the sample.
Require that an input point and an intersection point of two original
circles have different horizontal coordinates unless they are the same
point. Require also that an input point have a horizontal coordinate
different from that of a horizontal extremum of an original circle,
unless it is that extremum. Only finitely many directions are excluded.
For the latter assertion, if \(c\) is a circle center and \(p\) is a
point, an equality with an extremal abscissa has the form
\(e\cdot(p-c)=\pm1\) for a unit horizontal direction \(e\), and has
at most finitely many solutions in \(e\). This auxiliary direction
does not change the coordinates used elsewhere in the proof.

Ignore repetitions among the sampled circles. Express each sampled
circle as its upper and lower semicircle graphs on their open horizontal
domains. The event points are their horizontal extrema and all their
pairwise intersections, including tangencies. Between consecutive event
abscissas, the graphs have fixed vertical order; consecutive graphs,
with infinite outer boundaries allowed, bound open gaps.

Sweep these gaps from left to right. Continue a gap across an event
abscissa precisely when its two boundary graphs and their adjacency
continue unchanged and its closed limiting vertical segment contains no
event point. Otherwise start a new piece on the right. Each resulting
piece is the open region between two fixed semicircle graphs, or infinite
boundaries, over an open interval. It contains no part of a sampled
circle. In particular, a continued gap has no start, end, tangency or
crossing inside it at the abscissa across which it continues. An input
point at an event abscissa must, by the choice of direction, itself be an
event point and therefore lies on a sampled circle. Thus every point not
deleted is in one of the open pieces.

There are \(O(1+s^2)\) pieces, including in the presence of multiple
intersections or several events with the same abscissa. A new right-hand
gap can start only if an event point lies in its closed limiting vertical
segment. If \(c_x\) distinct sampled circles pass through an event point
\(x\), only \(O(1+c_x)\) right-hand gaps can have this property, by
their vertical ordering. Summing over event points costs \(O(s+s^2)\):
extrema contribute \(O(s)\), and at an intersection with \(c_x\ge2\)
one may charge \(c_x\) to \(2\binom{c_x}{2}\), while each pair of
circles has at most two intersections. This also covers extrema that
coincide with intersections. There is only a constant number of initial
unbounded gaps.

Finally, all possible pieces have a description from the full circle
set using a bounded number of circles. Each of their two horizontal
endpoints is an extremal or intersection abscissa, or infinity; each
boundary graph is specified by a circle and a sign, or is infinite.
Consequently at most \(C_0(1+l)^6\) candidate regions can occur, with
an absolute constant \(C_0\). For any fixed candidate region intersecting
more than \(l/R\) circles, the probability that no sampled circle
intersects it is at most
\[
 (1-1/R)^s\le e^{-s/R},
\]
with the evident interpretation when \(R=1\). An actual open piece
avoids every sampled circle. A union bound over the candidate regions,
with \(C\) sufficiently large, therefore gives the asserted conflict
bound with failure probability at most \(1/10\).
\end{proof}

\begin{proposition}\label{inc:incidence}
For \(m\) distinct points and \(l\) distinct unit circles in the plane,
\begin{equation}\label{inc:incidence-bound}
 I\le C\log^2(2+m+l)
       \bigl(m^{2/3}l^{2/3}+m+l\bigr),
\end{equation}
where \(C\) is an absolute constant. In particular, when \(m,l\) are
bounded by a fixed power of \(t\), the logarithmic factor is
\((\log t)^{O(1)}\).
\end{proposition}

\begin{proof}
Sz\'ekely's point--curve incidence theorem
\cite[Theorem~8]{Szekely1997} gives the stronger bound
\(I=O(m^{2/3}l^{2/3}+m+l)\) when distinct curves meet in at most two
points and distinct points belong to at most two curves. Distinct unit
circles satisfy both hypotheses. If simple curves are taken to be arcs,
cut each circle at two points outside the finite input point set; the
resulting \(2l\) arcs satisfy the same pair bounds and preserve every
incidence. The logarithmically weakened form
\eqref{inc:incidence-bound} follows, including the empty-set cases.
\end{proof}

\subsection{Degree control and a maximizing rectangle}

We now use the cutting to obtain the prescribed imbalance, and then
choose a rectangle of the incidence graph that cannot gain too much
edge density by further restriction. This last maximization supplies
both expansion and lower degree bounds.

\begin{proposition}\label{inc:piece}
From a sequence of \(n=t^3\) real planar points with at least
\(t^{4-o(1)}\) ordered unit pairs, one can obtain real algebraic point
and center sets \(P,Q\), and a nonempty graph \(E\subseteq P\times Q\)
consisting of unit incidences, with the following properties. There is
a parameter \(a=a(t)>0\) such that \(a\to0\),
\(a\log t\to\infty\), and
\begin{equation}\label{inc:scales}
 |P|=t^{1+2a+o(a)},\qquad
 |Q|=t^{2+a+o(a)},\qquad
 |E|=t^{2+2a+o(a)}.
\end{equation}
Writing \(d(p),d(q)\) for graph degrees, uniformly over all vertices,
\begin{equation}\label{inc:degrees}
 d(p)=t^{1+o(a)},\qquad d(q)=t^{a+o(a)}.
\end{equation}
The uniform law on \(E\) has marginals
\(\mu(p)=d(p)/|E|\) and \(\kappa(q)=d(q)/|E|\) satisfying
\begin{equation}\label{inc:marginal-atoms}
 \mu(p)\ge\frac{\gamma}{|P|},\qquad
 \kappa(q)\ge\frac{\gamma}{|Q|},\qquad
 \gamma=\frac{27}{50}=0.54.
\end{equation}
Set \(\eta=\gamma-\tfrac12=1/25\). There is a single sequence
\(\xi_t\to0\) such that, for every \(F\subseteq P\) and
\(G\subseteq Q\), with counting fractions
\(u=|F|/|P|\) and \(v=|G|/|Q|\),
\begin{equation}\label{inc:expansion}
 \Prob_E(p\in F,q\in G)
 \le \sqrt{uv}\bigl((uv)^\eta+\xi_t\bigr).
\end{equation}
The sequence \(\xi_t\) is independent of the chosen subsets.
\end{proposition}

\begin{proof}
First use \Cref{inc:specialization}. Choose \(a\to0\) sufficiently
slowly that the initial incidence count is at least \(t^{4-o(a)}\)
and every fixed power of \(\log t\) is \(t^{o(a)}\). In particular
\(a\log t\to\infty\). Apply \Cref{inc:cutting} with
\(R=t^{1-a}\) and \(s=O(t^{1-a}\log t)\). There are \(t^3\)
points and circles. Incidences on the sampled circles number at most
\(s t^3\), and incidences at deleted points with other circles number
at most \(2s t^3\). Thus the deletion costs
\(O(t^{4-a}\log t)=o(t^{4-o(a)})\).

Each remaining region meets at most \(t^{2+a}\) circles. Further split
its assigned point set into sets of size at most \(2t^{1+2a}\).
The total number of resulting pieces is at most the original region
count plus \(O(t^3/t^{1+2a})\), hence at most
\(t^{2-2a+o(a)}\). Some piece has at least
\(t^{2+2a-o(a)}\) incidences. Denote its two vertex sets temporarily
by \(P_0,Q_0\), and its edge count by \(e_0\). At this stage
\[
 |P_0|\le 2t^{1+2a},\qquad |Q_0|\le t^{2+a},\qquad
 t^{2+2a-o(a)}\le e_0\le t^{2+2a+o(a)},
\]
where the upper bound follows from \eqref{inc:incidence-bound}.

Choose \(b=o(a)\) so slowly that \(b\log t\to\infty\) and \(b\)
dominates all exponent errors in these estimates and the logarithmic
factor in \eqref{inc:incidence-bound}. Delete point vertices of degree
greater than \(t^{1+b}\), and center vertices of degree greater than
\(t^{a+b}\), with degrees measured before these deletions. The counts
of such vertices are at most
\[
 t^{1+2a-b/2}\quad\hbox{and}\quad t^{2+a-b/2},
\]
respectively, for all sufficiently large \(t\). Apply
\eqref{inc:incidence-bound} to the excessive vertices on one side and
all vertices on the other. In either case the mixed term is at most
\(t^{2+2a-b/3+o(b)}\); the two linear terms are bounded by
\(t^{1+2a+o(b)}+t^{2+a+o(b)}\). Their sum is \(o(e_0)\).
Both deletions therefore retain \(t^{2+2a+o(a)}\) edges and give
the degree upper bounds \(t^{1+o(a)}\), \(t^{a+o(a)}\).
Dividing the retained edge count by these two upper bounds shows that
both retained vertex counts attain their stated upper scales within
\(t^{o(a)}\). Relabel this trimmed graph as \((P_0,Q_0,E_0)\).

For subsets of these trimmed vertex sets, let \(u,v\) be their
respective counting fractions and let \(w\) be their captured edge
fraction. The incidence bound and the maximum-degree bounds give,
with one uniform factor \(T_t=t^{o(a)}\),
\begin{align}
 w&\le T_t\bigl((uv)^{2/3}+t^{-a}v+t^{-1}u\bigr),
 \label{inc:preliminary-expansion}\\
 w&\le T_t\min(u,v).\nonumber
\end{align}
Taking the minimum with the degree bounds separately in the two linear
terms, and using \(\min(x,y)\le\sqrt{xy}\), yields
\begin{equation}\label{inc:preliminary-square-root}
 w\le T_t\sqrt{uv}
       \bigl((uv)^{1/6}+2t^{-a/2}\bigr),
\end{equation}
after increasing \(T_t\) if necessary. For example,
\(\min(t^{-a}v,u)\le t^{-a/2}\sqrt{uv}\); the other linear term
uses \(t^{-1/2}\le t^{-a/2}\) for large \(t\).

Put \(\tau=t^{-a/8}\). Among all induced rectangles with counting
fractions \(u,v\ge\tau\) and captured edge fraction \(w\), maximize
\[
 \frac{w}{(uv)^\gamma},\qquad \gamma=\frac{27}{50}.
\]
We choose \(\gamma>1/2\) so that a bound by \((uv)^\gamma\)
improves on \(\sqrt{uv}\) for small sets. The inequality
\(\gamma<2/3\), where \(2/3\) is the incidence exponent in
\eqref{inc:preliminary-expansion}, will ensure that the maximizing
rectangle retains the scales \eqref{inc:scales}.
This is a maximum over a nonempty finite family, and its value is at
least one because the full rectangle is feasible. Write \(U,V\)
for the maximizing counting fractions and \(W\) for its captured
edge fraction. We claim
\begin{equation}\label{inc:winning-fractions}
 U,V,W\ge t^{-o(a)}.
\end{equation}
Indeed \eqref{inc:preliminary-square-root} gives
\[
 1\le T_t\left((UV)^{2/3-\gamma}
       +2t^{-a/2}(UV)^{-\eta}\right).
\]
By feasibility, \(UV\ge t^{-a/4}\), so the second term on the right
is at most \(2t^{-49a/100+o(a)}\), which tends to zero. If
\(UV\le t^{-ca}\) along a subsequence for any fixed \(c>0\), the
first term tends to zero as well, since
\(2/3-\gamma=19/150>0\). This is impossible. Thus
\(UV\ge t^{-o(a)}\), which implies the same for each of \(U,V\),
and \(W\ge(UV)^\gamma\) gives the claim for \(W\).

Let \((P,Q,E)\) be this maximizing rectangle, retaining precisely its
edges from the trimmed graph. The scale estimates
\eqref{inc:scales} and the degree upper bounds survive by
\eqref{inc:winning-fractions}. For further subsets with fractions
\(u,v\) inside \(P,Q\), optimality gives
\[
 \Prob_E(F,G)\le(uv)^\gamma
\]
whenever \(Uu,Vv\ge\tau\). Otherwise
\(uv\le\tau/\min(U,V)=t^{-a/8+o(a)}\). Apply
\eqref{inc:preliminary-square-root} with the old fractions \(Uu,Vv\)
and divide by \(W\). Uniformly in all such infeasible subsets, this
gives
\[
 \Prob_E(F,G)
 \le t^{o(a)}\sqrt{uv}
       \bigl((uv)^{1/6}+2t^{-a/2}\bigr)
 \le\xi_t\sqrt{uv},
\]
where one may take \(\xi_t=t^{-a/48+o(a)}\to0\), with the exponent
error depending only on the chosen rectangle and the previous uniform
bounds. Combining the feasible and infeasible cases proves
\eqref{inc:expansion}. Empty subsets satisfy it trivially.

Finally \eqref{inc:winning-fractions} implies that removing any single
vertex of the maximizing rectangle remains feasible for all sufficiently
large \(t\). Optimality with a point \(p\) removed gives
\[
 1-\mu(p)\le(1-|P|^{-1})^\gamma\le1-\gamma|P|^{-1},
\]
the last inequality following from concavity. The center comparison is
identical. This proves \eqref{inc:marginal-atoms}, and multiplying by
\(|E|\) supplies the lower degree bounds in \eqref{inc:degrees}.
\end{proof}

In particular, the minimum degree on either side tends to infinity,
and \(\max_p\mu(p)=t^{-1+o(1)}\). The lower marginal bounds also
convert \eqref{inc:expansion} to the form needed for probability laws.
For \(h=\mu(F)\), \(k=\kappa(G)\), we have
\(|F|/|P|\le h/\gamma\), \(|G|/|Q|\le k/\gamma\), and hence
\begin{equation}\label{inc:marginal-expansion}
 \Prob_E(F,G)\le C\sqrt{hk}\bigl((hk)^\eta+\xi_t\bigr),
\end{equation}
with a fixed absolute \(C\). The same \(\xi_t\) works for all subsets.

\subsection{Matching labels and uniformly rare sets}

We will repeatedly assign finite labels to the vertices, with no bound
on the number of labels. Expansion says that if almost all edges match,
one label occupies almost all the mass. A second application of
expansion makes the error proportional to the actual mismatch, even
when that mismatch is smaller than \(\xi_t\).

\begin{lemma}\label{inc:matching}
For the sequence of edge laws in \Cref{inc:piece}, there are absolute
constants \(s_0>0,C<\infty\) and an index after which the following
holds uniformly for every pair of label maps
\(\ell_P:P\to\mathcal A\), \(\ell_Q:Q\to\mathcal A\) into a
finite set. If
\[
 s=\Prob_E\bigl(\ell_P(p)\ne\ell_Q(q)\bigr)\le s_0,
\]
some label \(a_0\in\mathcal A\) satisfies
\[
 \mu\{\ell_P\ne a_0\}+\kappa\{\ell_Q\ne a_0\}\le Cs.
\]
Moreover, for independent \(p,p'\sim\mu\), at every value of \(s\),
\begin{equation}\label{inc:point-label-mismatch}
 \Prob\bigl(\ell_P(p)\ne\ell_P(p')\bigr)\le Cs.
\end{equation}
The analogous statement holds for independent centers of law \(\kappa\).
\end{lemma}

\begin{proof}
Let \(u_i,v_i\) be the counting fractions of the two classes with label
\(i\), including zero fractions. Sum \eqref{inc:expansion} over the
matched rectangles. Cauchy--Schwarz gives
\(\sum_i\sqrt{u_iv_i}\le1\), and therefore
\[
 1-s\le\max_i(u_iv_i)^\eta+\xi_t.
\]
For sufficiently small fixed \(s_0\) and sufficiently large \(t\),
one label \(a_0\) thus has both counting fractions arbitrarily close
to one. Apply \eqref{inc:expansion} to either complementary class
against the whole opposite side. The marginal exception masses
\[
 h=\mu\{\ell_P\ne a_0\},\qquad
 k=\kappa\{\ell_Q\ne a_0\}
\]
are then as small as any prescribed fixed constant, uniformly over
the label maps. This first conclusion need not be proportional to \(s\).

To obtain proportionality, put
\(J=\Prob_E(\ell_P\ne a_0,\ell_Q\ne a_0)\).
By \eqref{inc:marginal-expansion} and \(\sqrt{hk}\le(h+k)/2\),
we may choose the preceding constants so that
\(J\le(h+k)/4\). The events on which exactly one endpoint has label
\(a_0\) are mismatches, so
\[
 s\ge h+k-2J\ge\tfrac12(h+k).
\]
This proves the first assertion, including when \(s=0\), with no
additive error. Two independent point labels differ with probability
at most \(2h\), giving \eqref{inc:point-label-mismatch} for small
\(s\). For \(s>s_0\), the bound follows by increasing \(C\) so
that \(Cs_0\ge1\). The center assertion is identical.
\end{proof}

Define the \emph{wedge law} on triples \((p,q,r)\in P\times Q\times P\)
by first choosing \(q\sim\kappa\) and then independently choosing
\(p,r\) according to the conditional edge law given \(q\). Explicitly,
\begin{equation}\label{inc:wedge-law}
 \Prob(p,q,r)=
 \begin{cases}
 (|E|d(q))^{-1},&(p,q),(r,q)\in E,\\
 0,&\text{otherwise}.
 \end{cases}
\end{equation}
Each point marginal is \(\mu\), and each of the two point--center
marginals is the uniform edge law.

\begin{lemma}\label{inc:rare}
Let \(\delta_t\to0\) be any nonnegative sequence. The estimates below
hold uniformly over all indicated subsets, with error factors depending
only on \(\delta_t,\xi_t\) and the fixed constants in
\eqref{inc:marginal-expansion}:
\begin{enumerate}[label=\textup{(\roman*)}]
\item If \(F\subseteq P\), \(G\subseteq Q\), and
\(h=\mu(F)\le\delta_t\), while \(k=\kappa(G)\) is arbitrary, then
\begin{equation}\label{inc:one-small-set}
 \Prob_E(F,G)\le C(\delta_t^\eta+\xi_t)\sqrt{hk}
                 =o(1)\sqrt{hk}.
\end{equation}
The same assertion holds with the sides interchanged.
\item If also \(k\le\delta_t\), then
\begin{equation}\label{inc:joint-rare-edge}
 \Prob_E(F,G)\le\tfrac C2(\delta_t^{2\eta}+\xi_t)(h+k)
                 =o(1)(h+k).
\end{equation}
\item If \(F\subseteq P\) and \(h=\mu(F)\le\delta_t\), then
under the wedge law,
\begin{equation}\label{inc:joint-rare-wedge}
 \Prob(p\in F,r\in F)=o(1)h.
\end{equation}
\end{enumerate}
In particular, the error factors may be taken outside sums or integrals
over any family of sets whose marginal masses obey the same uniform
bound \(\delta_t\).
\end{lemma}

\begin{proof}
The first two assertions follow directly from
\eqref{inc:marginal-expansion}, using \(hk\le\delta_t\) in the
first case and \(hk\le\delta_t^2\) and
\(\sqrt{hk}\le(h+k)/2\) in the second.

For the third, let \(g(q)=\Prob(p\in F\mid q)\), so that
\(\E_\kappa g=h\) and the probability in question is
\(\E_\kappa g^2\). There is nothing to prove when \(h=0\).
For any \(d\in(0,1]\), put \(G_d=\{q:g(q)>d\}\). Markov's
inequality gives \(\kappa(G_d)\le h/d\). Splitting at this threshold,
\[
 \E_\kappa g^2
 \le dh+\E_\kappa[g\1_{G_d}]
 =dh+\Prob_E(F,G_d)
 \le h\left[d+C d^{-1/2}
                  \bigl((\delta_t^2/d)^\eta+\xi_t\bigr)\right].
\]
Choose \(d=d_t\to0\) sufficiently slowly that
\(\delta_t^{2\eta}d_t^{-1/2-\eta}\to0\) and
\(\xi_t d_t^{-1/2}\to0\). The bracket then tends to zero,
independently of \(F\). This proves the uniform assertion and hence
its integrated form.
\end{proof}

\subsection{The information carried by a wedge}

The small point side and the two-center codegree bound together ensure
that the two point endpoints of a wedge share little information. We
record the atom estimates as well, since they later convert a probability
of wedges into a count of distinct point pairs.

For probability laws \(a,b\) on the same finite set, define their relative
entropy by
\begin{equation}\label{inc:relative-entropy}
 \KL(a\Vert b)=\sum_x a(x)\log\frac{a(x)}{b(x)},
\end{equation}
with a zero summand when \(a(x)=0\) and value \(+\infty\) if
\(a(x)>0=b(x)\).

\begin{lemma}\label{inc:entropy}
For the wedge law \eqref{inc:wedge-law}, let
\(d_{\min,Q}=\min_{q\in Q}d(q)\). If \(p\ne r\), then
\begin{equation}\label{inc:wedge-atoms}
 \Prob(p,r)\le\frac{2}{|E|d_{\min,Q}}
                 =t^{-2-3a+o(a)}.
\end{equation}
The diagonal probability is
\begin{equation}\label{inc:wedge-diagonal}
 \Prob(p=r)=\frac{|Q|}{|E|}=t^{-a+o(a)}=o(1).
\end{equation}
The mutual information of the point endpoints satisfies
\begin{equation}\label{inc:information-bound}
 D:=\KL\bigl(\mathcal L(p,r)\,\Vert\,\mu\otimes\mu\bigr)
   =o(\log t).
\end{equation}
\end{lemma}

\begin{proof}
For fixed \(p,r\), summing \eqref{inc:wedge-law} gives
\[
 \Prob(p,r)=\frac1{|E|}
       \sum_{q:(p,q),(r,q)\in E}\frac1{d(q)}.
\]
There are at most two summands when \(p\ne r\), proving the
off-diagonal estimate by \eqref{inc:scales}--\eqref{inc:degrees}.
Summing instead over the diagonal yields
\[
 \Prob(p=r)=\frac1{|E|}\sum_{q\in Q}
                 \frac{d(q)}{d(q)}=\frac{|Q|}{|E|}.
\]
The asymptotics follow from the same estimates and
\(a\log t\to\infty\).

For a finite random variable \(X\), write
\(\mathsf H(X)=-\sum_x\Prob(X=x)\log\Prob(X=x)\), with
\(0\log0=0\). Every off-diagonal atom is at most
\(M_t=2/(|E|d_{\min,Q})\). Dropping the nonnegative diagonal
entropy contributions therefore gives
\[
 \mathsf H(p,r)
 \ge\bigl(1-\Prob(p=r)\bigr)\log(1/M_t)
 \ge(2-o(1))\log t.
\]
On the other hand,
\(\mathsf H(p)=\mathsf H(r)\le\log|P|
=(1+o(1))\log t\). Expanding the definition of relative entropy
gives
\[
 0\le D=\mathsf H(p)+\mathsf H(r)-\mathsf H(p,r)
       =o(\log t),
\]
as required.
\end{proof}

We retain this graph and these probability laws for the rest of the
proof: the uniform edge law, its point and center marginals, and the
wedge law. The expansion consequences hold for every finite labeling
and every sufficiently small exceptional set, with constants independent
of any tests subsequently assigned to the vertices.
